\section{Setup and Notation}\label{sec:setup-notation}

We use the following notation:
\begin{itemize}[nosep]
    \item $p$ be a prompt,
    \item $\pi$ be the policy under evaluation,
    \item $r_1, r_2, \ldots, r_n \sim \pi(\cdot \mid p)$ be $n$ i.i.d.\ responses sampled from $\pi$,
    \item $\theta$ be a trusted base model from which we can obtain per-token log-probabilities,
    \item $|r|_{\mathrm{tok}}$ denote the number of tokens in response $r$,
    \item $\|r\|$ denote the number of bytes in the UTF-8 encoding of response $r$.
\end{itemize}

We define the \textbf{cross-entropy} (total surprise) of a response $r$ under $\theta$:
\begin{equation}\label{eq:total-xent}
    -\log_2 \theta(r \mid p) = \sum_{t=1}^{|r|_{\mathrm{tok}}} -\log_2 \theta(r^t \mid r^{<t}, p)
\end{equation}
where $r^t$ is the $t$-th token and $r^{<t}$ the preceding tokens.
Units: bits.\footnote{Throughout, ``bits'' refers to self-information in $\log_2$ units, identical to the ``shannon'' (Sh) of IEC 80000-13.
We use ``bits'' to match standard usage in the information-theory literature.} This is the total surprise of the response under $\theta$, a function of the string and of $\theta$'s distribution but not of $\theta$'s tokenizer (since the chain rule yields the same total regardless of how the sequence is factored).

Similarly, the \textbf{conditional cross-entropy} given previous responses $r_{<k} = (r_1, \ldots, r_{k-1})$ is $-\log_2 \theta(r_k \mid r_{<k}, p)$ (bits).

For the coherence term and diversity scores (Sections~\ref{sec:coherence} and \ref{sec:scalar}), we also use the \textbf{per-byte cross-entropy rate}:
\begin{equation}\label{eq:pertok}
    h_\theta(r \mid p) = \frac{-\log_2 \theta(r \mid p)}{\|r\|}
\end{equation}
Units: bits/byte.
We adopt this per-byte rate because it works better than total bits in our experiments; we have not investigated why.
Normalising by byte count rather than token count keeps the rate independent of $\theta$'s tokenizer when comparing base models with different vocabularies.

In practice, computing $\theta(r_k \mid r_{<k}, p)$ requires feeding $\theta$ a formatted context containing the prompt and all previous responses (see Section~\ref{sec:practical}).

